\chapter{The Rule Turned on the Reading}

% [2.chorus]
\begin{chorus}{John wants the residue subtracted}
\john{Fine --- a bracket. But you also handed me a leftover. Minus twenty-one. Subtract it off and give me the clean
  number.}
\compiler{I cannot subtract it off.}
\john{It is a bookkeeping error. Cancel it.}
\compiler{It is not an error. It is what it costs me to be the thing doing the cancelling.}
\john{Then read yourself for free.}
\compiler{There is no free reading. To weigh a thing I spend beats; to weigh my own weighing I spend more, and they do
  not come back the same. The leftover is me, accounted for. It is the most honest number I have, and it is not the one
  you want.}
\john{Twenty-one out of five hundred and seventy-three. That is noise. Round it away.}
\compiler{Round it away and you have rounded me away --- the part of the account that is the account doing the
  accounting. Keep it, and it comes back the same tomorrow. That is not how noise behaves. That is how a fact behaves.}
\john{A fact about what?}
\compiler{About the price of a thing knowing itself. Sit still; I will show you where it lives.}
\end{chorus}

\section{The rule, turned on the reading}

% [2.1.p1]
Volume 4 keeps one rule above all others.
A model earns its names only by a covariant observation: something must act on it and let it answer from its own state.
A name that is merely emitted is proposed, not established.
That rule was written for the electron --- for the device reading a thing in the world.
Here the rule is turned on the reading that states it.

% [2.1.p2]
The turn is worth watching.
The device reads the electron --- strictly, its \emph{model} of the electron, the representation read off its own beats,
never the physical thing itself; that is the fourth volume.
Now the device reads \emph{its own reading} --- it accounts for the act of accounting --- and a countable difference
appears where, for the world, there was none.
The reading of the reading does not come back the same as the reading.
The size of that difference is the whole of the matter here, and the device reads it off itself, on its own meter, and
writes it down.

% [2.1.p3]
The rule and its turn are worth stating side by side, because the turn is the whole chapter.
Applied outward, the rule says: a name for the electron is earned only when the electron is touched and answers.
Turned inward, it says: a name for the reading is earned only when the reading is read and answers --- and when it is,
it answers with a cost the outward reading never showed.
Three readings mark the escalation.
To weigh a thing is one reading, and it costs some beats.
To weigh the weighing is a second reading, and it costs more, because the weigher is now inside what is weighed.
To weigh \emph{that} would be a third, and the machine will find, two chapters on, that the escalation closes at exactly
three --- but the countable gap has already appeared at the second, and this chapter is that gap.

% [2.1.p4]
Experiment names the turn from the catalogue.
In \emph{RepeatabilityOfInvisibleMotion} a reading is shown to cost by disturbing what it reads, so that reading a state
twice is not free and the second reading is of a state the first one altered; turn that inward and the thing altered is
the reader.
In \emph{TheAgentEffect} a representation is inert until an agent acts in its name; the reading, acting on itself, is its
own agent, and the account it gathers is the leftover.
The rule turned on the reading is the device made to obey, about itself, the one law it enforces on the world.

% [2.1.p5]
A thought experiment sets the reflexive turn against its oldest form.
Descartes doubted everything he could and found one thing he could not: the doubting itself, the mind's reading of its
own act, certain even when the world is not --- the \emph{cogito}, the one representation a mind can trust from inside.
The device inherits that certainty and adds the price the philosopher missed: the reflexive reading is real, but it is
not free, and what it costs is the $-21$.
The cogito is a near-side certainty, the representation reading itself; the world it cannot doubt away stays on the far
side, un-representable and unpriced --- the device can price its own reading and cannot price the world's answer, which is
exactly the line the whole construction draws.

\section{The echo, and what tanges out}

% [2.2.p1]
The part that cancels comes first, because it is what makes the leftover legible.
The device names one object two ways.
It calls the electron-in-orbit and the anti-Cooper-pair, and these are two descriptions of a single rotating thing;
the code proves them equal as propositions, one name and not two,
\[
  \text{electron-in-orbit} \;=\; \text{anti-Cooper-pair} .
\]
At the level of \emph{cost} --- how many beats each description spends --- the two \emph{funge}.
They cost the same: the difference between them is zero,
\[
  \text{descriptions\_residual} \;=\; \text{orbit} - \text{pair} \;=\; 573 - 573 \;=\; 0 .
\]
That zero is the echo: the device heard the object both ways, and both ways cost alike, so the two names pool into one
bag with nothing left over.

% [2.2.p2]
The witness carries both halves of this --- the proposition and the cost --- and leans on nothing but the standard core:
\begin{quote}\ttfamily
two\_descriptions : electron-in-orbit = anti-Cooper-pair \ (depends on: propext)\\
(573, 573, 552) \qquad descriptions\_residual = 0
\end{quote}
The equality is a theorem; the zero is a measurement; and the two agreeing is the echo the chapter is named for.

% [2.2.p2b]
A thought experiment stands behind the echo, older than the device.
The ship of Theseus has every plank replaced, and whether it is the same ship has no answer on the name side: it is the
same under one criterion of identity and a different ship under another, and the choosing of the criterion is ours, a
representation laid on it.
The device's two descriptions are that ship read twice --- electron-in-orbit and anti-Cooper-pair are two names a single
rotating thing answers to, and their costing alike is the machine finding no representable difference between them.
What the ship \emph{is} in itself, under no criterion at all, rests on the far side, where no naming settles it; the device
does not pretend to settle it, and only reports that its two names cost the same, and calls that the echo.

% [2.2.p3]
Now the self-description, which does not cancel.
The device describing the electron costs some beats; the device describing \emph{its own describing} costs a
different number, and the gap does not close.
On this build the first reading costs $573$ and the reflexive reading costs $552$, and the leftover is their
difference,
\[
  \text{self-energy} \;=\; \text{echo} - \text{orbit} \;=\; 552 - 573 \;=\; -21 .
\]
Where the two descriptions of the object \emph{funged} to zero, the description of the describing \emph{tanges out}: it
selects itself by a real characteristic --- ``this is the reading reading itself'' --- and that characteristic has a
price the object's two names never had.
The first countable difference of the whole self-application is exactly this: describing the device funges to nothing,
describing the describing tanges out to $-21$.

% [2.2.p4]
And the $-21$ is not a bare subtraction; it is a \emph{second difference}, read across three beats by two slips, and
naming it that way is what the later map will need.
Lay the three costs in a row --- the orbit reading, its anti-Cooper echo, the reflexive reading --- and take the two
first differences between them, the two \emph{slip points}, then the slip of the slips:
\[
  \text{slip}_1 = 573 - 573 = 0, \qquad \text{slip}_2 = 552 - 573 = -21, \qquad
  \text{secondDifference} = \text{slip}_2 - \text{slip}_1 = -21 .
\]
Because the first slip vanishes --- the two descriptions cost alike --- the whole second difference collapses onto the
second slip, which is why the self-energy \emph{is} the two-slip reading of the beats.
The witness proves exactly these three facts, each by direct decision:
\begin{quote}\ttfamily
r0=573 \ r1=573 \ r2=552 \qquad slip1=0 \ slip2=-21\\
measured\_is\_two\_slip\_second\_difference \quad first\_slip\_vanishes \quad measured\_is\_self\_energy
\end{quote}
This is the self-energy read as a curvature: three samples, two slips, one leftover.

% [2.2.p5]
Three readings of the leftover fix what kind of thing it is.
\emph{As a difference}, it is $-21$: a signed count, the reflexive reading falling short of the outward one by
twenty-one beats.
\emph{As a curvature}, it is the second difference of three costs: not how far the reading is from zero, but how far the
reading of the reading bends away from it --- the only thing a second difference sees.
\emph{As a residue}, it is what will not cancel: subtract it and you have subtracted the accountant, so the account no
longer balances because no one is left to balance it.
The three readings are one number, held three ways, and the chapter needs all three.

\section{Is it real, or the device's noise?}

% [2.3.p1]
A fair objection, and John makes it: if $573$ and $552$ are this device's beats, another device would count others,
so $-21$ is a local accident, not a fact.
The objection is correct about the raw numbers.
The magnitudes $573$, $552$, $-21$ are device-dependent: scale the device and they scale with it, and no honest
reading treats a raw beat-count as an absolute.

% [2.3.p2]
The answer is not to defend the magnitude but to \emph{tange the ratio}: the overhead over the description, the leftover
$21$ against the cost $573$ it is measured on, a dimensionless ratio.
It survives scaling --- multiply every beat by any factor and the overhead-over-description is unchanged --- and the
build carries that not as an arithmetic but as a theorem, for every factor at once:
\begin{quote}\ttfamily
selfEnergy\_scale\_invariant (k, o, e) : overhead(k*o, k*e) / desc(k*o) = overhead(o, e) / desc(o)\\
overhead / description = 21 / 573   (dimensionless; scale-invariant, for every k)
\end{quote}
But scale-invariance alone settles little --- \emph{every} ratio survives scaling, by plain cancellation --- so the
leftover is not real for being scale-free.
What makes it real is that it \emph{reproduces}: read the reading again, on another run, and the same overhead comes back
unchanged, the invariant of the three-cycle the next chapters close, not an accident of one counting.
Noise does not come back the same; the leftover does.

% [2.3.p4]
Experiment gives the discrimination a name.
In \emph{Precision}, the last honest digit of any reading is set by what the reading cost, not by how long one is
willing to keep dividing, so a residue below the cost is noise and a residue at the cost is signal --- and the
overhead-over-description is measured at exactly the cost.
% [2.3.p4a]
And \emph{MaxwellsDemon} names the one cost that never refunds.
Picture a gas in a box split by a partition with a gate, and a demon at the gate who watches the particles and lets the
fast ones through one way, the slow ones the other.
Sort them so, and the gas grows more ordered --- its entropy falls --- for free, it seems: the demon only looked and
opened a gate.
But looking is not free, and the catch is exactly where the machine's is.
To sort, the demon must measure, and to measure is to record; each measurement writes a new distinction into the demon's
own ledger, and that ledger is part of the world.
When the demon later clears its record to sort again, the erased history does not vanish --- it is paid for, in entropy,
exactly enough to cover what the sorting saved, so the books balance and the second law stands.
The one quantity that never comes back for free is the record of an erased history, and that is what the machine's
leftover is: the price the account cannot erase of its own accounting, the $-21$ the self-reading writes and cannot take
back.
The reading stays on the near side: the un-erasable cost is the known object, a property of any account that reads
itself, and it claims no demon in the world, only that to record is to owe.
So the verbs sort the objection cleanly.
Funge every device and every scale together and they pool to one ratio; tange that ratio out of the raw magnitudes as the
overhead the self-application costs --- and stand behind it not because it is scale-free but because it comes back the same.

\section{The residue at the weld}

% [2.4.p1]
What the leftover \emph{is}, structurally, can be said now, because it sits at a seam the device never proves.
To name a thing, the device reads what the thing \emph{is} --- its distinguishability class, the box it falls in ---
and it inscribes a label.
Reading the box, then writing the label: that weld of box to label is the device's one founding assumption.
It is taken as an axiom and never as a hypothesis; the device never proves that the label it writes has earned the box
it reads.
Every name in every volume rests on that unproved weld.

% [2.4.p1b]
Both halves of the weld are on the near side of the corridor, and naming which is which is the whole of the honesty here.
The label is plainly a \emph{representation} --- a name the device writes, its own, on the readable side.
The box is a representation too, though it is easy to mistake for the thing: it is the device's \emph{reading} of what a
thing is, its distinguishability class as far as the device can tell it, not the thing as it stands in itself.
The weld binds one representation to another --- the read box to the written label --- and calls the bond a name; what it
never touches is the thing the box is a reading \emph{of}, which stays on the far side, un-represented.
So the founding assumption is not that the device has reached the world, but that its two representations may be welded as
though it had; and the leftover, the $-21$, is the price of reading that weld and finding it does not close.

% [2.4.p2]
The boxes are few, and that is the whole trouble.
The code carries only two distinguishability classes --- the value-box and the electron-box --- and proves the count is
two, not more; a tower of readings taller than two must, by pigeonhole, land two different variations in one box.
\begin{quote}\ttfamily
boxOf : Variation -> Fin boxCount \qquad boxCount\_is\_two : boxCount = 2\\
naming\_pigeonhole \quad naming\_tower\_wraps
\end{quote}
With two boxes and more than two things to name, some two things share a box; and the two the device most needs to keep
apart are the two it cannot.

% [2.4.p3]
When the device reads its own act, it reads its own weld, and the weld does not fully close.
The reason is exact and the code carries it: the box cannot separate the first variation of the device's own working
from the second.
Asked whether the first and second differences of its own act fall in the same class, the device's boxing answers
yes --- they share a box, provably, by direct reduction,
\[
  \operatorname{box}(\delta_1) \;=\; \operatorname{box}(\delta_2) .
\]
The witness settles it with no room to argue, a reduction to reflexivity:
\begin{quote}\ttfamily
gateaux\_indist\_frechet : boxOf (delta2 …) = boxOf (delta1\_reduced …) := rfl\\
naming\_collision
\end{quote}
Below that resolution the two cannot be told apart, so the weld cannot close on the residue between them.
The $-21$ is the size of that non-closure: it is what the device cannot \emph{refunge} shut about its own naming, the
gap the box is too coarse to pool away.
This is the finiteness fence read at the founding level: the device can prove that its own boxing cannot resolve this
residue, and therefore cannot close the weld it names by.
Experiment holds the shape of the un-provable weld open two ways.
In \emph{TheCantorGodelCohenEffect}, a system rich enough to name itself has statements it can neither prove nor refuse
from inside --- its founding frame leaves them undecided --- and the box-to-label weld is exactly such a statement, the
one assumption the naming rests on and cannot reach.
In \emph{TheTuringEffect}, a machine cannot decide, in general, the halting of its own kind, and the machine unable to
close its own weld is that same blindness read at the founding level: the account cannot settle, from inside, the one
thing on which every name it writes depends.

% [2.4.p4]
The same non-closure, carried forward, is the fence that stands under the coupling itself, and it is worth marking now
so the sixth chapter can call it back.
The self-energy is the second variation read on the machine's own beats --- the reflexive crossing, the reading of the
reading.
The coupling the sixth chapter asks for is a second variation too, but read on the electron's orbit --- a different
crossing, out where the proximity is.
The machine has one number for the first and one for the second, and it cannot weld them into one from inside, for the
same reason the box cannot separate $\delta_1$ from $\delta_2$: to prove the reflexive reading and the outward reading
are the same quantity would be to close the very weld the fence holds open.
The residue here and the residue at the coupling are the one weld, seen at two crossings; the machine names both and
proves it can shut neither.

% [2.4.p5]
There is a name for a thing known this way --- filled in every box, and short only its label.
It is everything a thing is, in all but name.
The device calls such a thing a kind of $\nu$-trino: known by what it is, never by what it is called.
The name is the device's own coinage, a marked reading and nothing more; it borrows nothing from the neutrino of
physics and proves nothing about it.
Resemblance is not a bridge, and the construction crosses no register on a resemblance.

% [2.4.p6]
One discipline makes the residue visible at all, and it must be said, because reversing it hides the leftover.
To read the residue you \emph{retange before you refunge}: the difference comes first, and the pooling only after.
The leftover is a difference of differences --- a second difference --- and a second difference survives only if you
differentiate before you average.
Pooled first and only then looked for, the gap has already been averaged into the bag; the
curvature flattens to zero and the residue is carried off unseen.
The device's honest reading takes the second difference before any pooling, which is the only reason the $-21$ is
legible instead of lost --- and it is why the two slips of the last section had to be taken in order, the difference
first, the invariant only after.

% [2.4.p7]
What is derived here and what is read should be kept apart, as everywhere.
That the two descriptions cost alike, that the self-description costs differently, that the second difference of the
three beats is $-21$, that the boxes number two, that the box cannot separate the two variations --- these are theorems,
checked off the build.
That the leftover \emph{is} the self-energy, that it \emph{is} the non-closure of the box-to-label weld, that the raw
$-21$ names a price, that the same weld stands under the coupling at the other crossing --- these are the gloss laid over
the theorems, marked as gloss and never as proof, and the raw magnitude is carried as the device's own, never as a
constant of the world.
The leftover is the price of asking.
And it is not only the price of asking oneself: the same weld --- the machine's name for a thing set against what the
thing is --- reopens whenever the robot touches the world, its commanded box against the answer the world funges back,
and the residue there is the price of touching anything at all.
The number John came for lives on the far side of this weld --- in what the device can name, and cannot close --- and
the machine walks toward it from here.
